\clearpage
\appendix
\section{SYCL Matrix Multiplication Example}
\label{appendix:sycl-matmul}
 
This appendix provides a \acrshort{sycl} implementation of a matrix multiplication kernel to illustrate the GPU programming workflow described in Section~\ref{subsubsec:gpu-programmation}. The code demonstrates how to orchestrate parallel computation on a device and includes annotations for each phase of the execution.

\begin{enumerate}
    \item \textbf{Matrix Dimension Initialization}
    Defining the matrix dimensions and calculating the required memory sizes in bytes.
\begin{lstlisting}
int M = 512, K = 512, N = 512;
size_t sizeA = M * K * sizeof(float);
size_t sizeB = K * N * sizeof(float);
size_t sizeC = M * N * sizeof(float);
\end{lstlisting}

    \item \textbf{Host Memory Allocation and Initialization}
    Allocating memory on the CPU (host) and populating the input matrices with initial data.
\begin{lstlisting}
// Allocate host memory
float* h_A = static_cast<float*>(malloc(sizeA));
float* h_B = static_cast<float*>(malloc(sizeB));
float* h_C = static_cast<float*>(malloc(sizeC));

// Initialize matrices with dummy data
for (int i = 0; i < M * K; i++) h_A[i] = 1.0f;
for (int i = 0; i < K * N; i++) h_B[i] = 2.0f;
\end{lstlisting}

    \item \textbf{SYCL Queue Creation}
    The queue serves as the primary interface for submitting command groups to the accelerator device.
\begin{lstlisting}
sycl::queue q;

// Optional - Output the selected device name
std::cout << "Running on: " 
          << q.get_device().get_info<sycl::info::device::name>() 
          << std::endl;
\end{lstlisting}

    \item \textbf{Device Allocation and Data Transfer}
    Allocating memory directly on the GPU (device) and transferring data from the host.
\begin{lstlisting}
// Allocate device memory using Unified Shared Memory (USM) device pointers
float* d_A = sycl::malloc_device<float>(M * K, q);
float* d_B = sycl::malloc_device<float>(K * N, q);
float* d_C = sycl::malloc_device<float>(M * N, q);

// Explicitly copy data to device memory
q.memcpy(d_A, h_A, sizeA);
q.memcpy(d_B, h_B, sizeB).wait(); // Ensure transfer is complete before proceeding
\end{lstlisting}

    \item \textbf{Kernel Dispatch}
    SYCL offers several methods for launching kernels on the device:
    \begin{enumerate}
        \item \textbf{Anonymous Lambda Function:} A concise method capturing variables by value using \lstinline{[=]}.
\begin{lstlisting}
q.parallel_for(sycl::range<2>(M, N), [=](sycl::id<2> index) {
    int row = index[0];
    int col = index[1];
    float sum = 0.0f;
    for (int i = 0; i < K; ++i) {
        sum += d_A[row * K + i] * d_B[i * N + col];
    }
    d_C[row * N + col] = sum;
}).wait();
\end{lstlisting}

        \item \textbf{Named Lambda:} Wrapping the kernel in a command group handler (\lstinline{submit}). This is particularly useful for debugging and profiling tools.
\begin{lstlisting}
q.submit([&](sycl::handler& h) {
    h.parallel_for<class MatMul>(sycl::range<2>(M, N),
                                 [=](sycl::id<2> index) {
        int row = index[0];
        int col = index[1];
        float sum = 0.0f;
        for (int i = 0; i < K; ++i) {
            sum += d_A[row * K + i] * d_B[i * N + col];
        }
        d_C[row * N + col] = sum;
    });
}).wait();
\end{lstlisting}

        \item \textbf{Functor Class:} Defining a function object to encapsulate the kernel logic, providing a clean separation between host and device code.
\begin{lstlisting}
struct MatrixMulFunctor {
    float *A, *B, *C;
    int M, K, N;

    // Constructor to capture host-side parameters
    MatrixMulFunctor(float* a, float* b, float* c,
                     int m, int k, int n)
        : A(a), B(b), C(c), M(m), K(k), N(n) {}

    // The call operator defines the kernel entry point
    void operator()(sycl::id<2> index) const {
        int row = index[0];
        int col = index[1];
        if (row < M && col < N) {
            float sum = 0.0f;
            for (int i = 0; i < K; ++i) {
                sum += A[row * K + i] * B[i * N + col];
            }
            C[row * N + col] = sum;
        }
    }
};

// Dispatching the kernel using the functor instance
q.submit([&](sycl::handler& h) {
    MatrixMulFunctor functor(d_A, d_B, d_C, M, K, N);
    h.parallel_for(sycl::range<2>(M, N), functor);
}).wait();
\end{lstlisting}
    \end{enumerate}

    \item \textbf{Result Retrieval}
    Copying the computed results from device memory back to the host for verification.
\begin{lstlisting}
q.memcpy(h_C, d_C, sizeC).wait();
std::cout << "Success! Element [0,0] is: " << h_C[0]
          << " (Expected: " << K * 2.0f << ")"
          << std::endl;
\end{lstlisting}

    \item \textbf{Memory Deallocation}
    Releasing both device and host memory to prevent leaks.
\begin{lstlisting}
sycl::free(d_A, q);
sycl::free(d_B, q);
sycl::free(d_C, q);
free(h_A);
free(h_B);
free(h_C);
\end{lstlisting}
    
\end{enumerate}